% GENERATED FILE. Edit canonical lessons, then run npm run generate:books.
% canonical-source-hash: fnv1a64:accff1ceab8e2452
% canonical-lessons: GU-C527-akhado, GU-C527-pahelvan, GU-C527-hariph, GU-C527-pachadvu, GU-C527-hamphavvu

\chapter{In the Wrestling Ring}
\label{ch:gu-in-the-wrestling-ring}
\hlchaptermodality{527}

\begin{chapteropening}
\textbf{By the end of this chapter:} \emph{I can describe a wrestling bout.}

Complete the last lesson of chapter 527: I can describe a wrestling bout.
\end{chapteropening}

\section[\texorpdfstring{\gu{અખાડો} (akhāḍo)}{અખાડો (akhāḍo)}]{\gu{અખાડો} (akhāḍo) — a wrestling ground}

\label{lesson:GU-C527-akhado}



\begin{quote}
Before the new one: say the Gujarati for a raga, then the Gujarati for to strike up.
\end{quote}

\subsection*{You'll want to know: \gu{અખાડો}}
\textbf{\gu{અખાડો}} — \emph{akhāḍo} — \textquotedblleft{}a wrestling ground\textquotedblright{}.

\begin{quote}
\gu{અમે} \gu{રોજ} \gu{સવારે} \gu{ગામના} \gu{અખાડામાં} \gu{કસરત} \gu{કરીએ} \gu{છીએ}.
\end{quote}

\begin{quote}
\emph{ame roj savāre gāmnā akhāḍāmā̃ kasrat karīe chīe.}
\end{quote}

\begin{quote}
Every morning we exercise at the village wrestling ground.
\end{quote}

A masculine noun.

\begin{cousinweb}
\textbf{\gu{અખાડો}} (\emph{akhāḍo}) is inherited from Sanskrit \emph{akṣavāṭa}, \textquotedblleft{}a wrestling ground\textquotedblright{}, and CDIAL lists the Gujarati word (CDIAL 39).
\end{cousinweb}

\begin{grammarlens}[title={What you've built}]
One more word for describing what you lived and felt.
\end{grammarlens}

\subsection*{Your turn}
\begin{itemize}
\raggedright
  \item \emph{Say it:} \emph{akhāḍo}
  \item \emph{Say it:} \emph{akhāḍo}, once more
  \item \emph{Say it:} \emph{ame roj savāre gāmnā akhāḍāmā̃ kasrat karīe chīe.} — Every morning we exercise at the village wrestling ground.
  \item \emph{Recall:} say the Gujarati for a raga, then the Gujarati for to strike up, then say \emph{akhāḍo} again
\end{itemize}

\begin{tcolorbox}[breakable,colback=teal!4,colframe=teal!35!black,title={Before you move on}]
What does \gu{અખાડો} mean? (\textquotedblleft{}A wrestling ground\textquotedblright{}.) Say it once more.
\end{tcolorbox}
\section[\texorpdfstring{\gu{પહેલવાન} (pahelvān)}{પહેલવાન (pahelvān)}]{\gu{પહેલવાન} (pahelvān) — a wrestler}

\label{lesson:GU-C527-pahelvan}



\begin{quote}
Before the new one: say the Gujarati for to strike up, then the Gujarati for a wrestling ground.
\end{quote}

\subsection*{You'll want to know: \gu{પહેલવાન}}
\textbf{\gu{પહેલવાન}} — \emph{pahelvān} — \textquotedblleft{}a wrestler\textquotedblright{}.

\begin{quote}
\gu{આ} \gu{પહેલવાન} \gu{ઘણો} \gu{મજબૂત} \gu{છે}.
\end{quote}

\begin{quote}
\emph{ā pahelvān ghaṇo majbūt che.}
\end{quote}

\begin{quote}
This wrestler is very strong.
\end{quote}

A masculine noun.

\begin{cousinweb}
\textbf{\gu{પહેલવાન}} (\emph{pahelvān}) is usually traced to Persian \emph{pahlavān}, \textquotedblleft{}a champion, a wrestler\textquotedblright{}.
\end{cousinweb}

\begin{grammarlens}[title={What you've built}]
One more word for describing what you lived and felt.
\end{grammarlens}

\subsection*{Your turn}
\begin{itemize}
\raggedright
  \item \emph{Say it:} \emph{pahelvān}
  \item \emph{Say it:} \emph{pahelvān}, once more
  \item \emph{Say it:} \emph{ā pahelvān ghaṇo majbūt che.} — This wrestler is very strong.
  \item \emph{Recall:} say the Gujarati for to strike up, then the Gujarati for a wrestling ground, then say \emph{pahelvān} again
\end{itemize}

\begin{tcolorbox}[breakable,colback=teal!4,colframe=teal!35!black,title={Before you move on}]
What does \gu{પહેલવાન} mean? (\textquotedblleft{}A wrestler\textquotedblright{}.) Say it once more.
\end{tcolorbox}
\section[\texorpdfstring{\gu{હરીફ} (harīph)}{હરીફ (harīph)}]{\gu{હરીફ} (harīph) — a rival}

\label{lesson:GU-C527-hariph}



\begin{quote}
Before the new one: say the Gujarati for a wrestling ground, then the Gujarati for a wrestler.
\end{quote}

\subsection*{You'll want to know: \gu{હરીફ}}
\textbf{\gu{હરીફ}} — \emph{harīph} — \textquotedblleft{}a rival\textquotedblright{}.

\begin{quote}
\gu{આજે} \gu{પહેલવાનનો} \gu{હરીફ} \gu{તેના} \gu{કરતાં} \gu{મોટો} \gu{છે}.
\end{quote}

\begin{quote}
\emph{āje pahelvānno harīph tenā kartā̃ moṭo che.}
\end{quote}

\begin{quote}
Today the wrestler's rival is bigger than he is.
\end{quote}

A masculine or feminine noun, by the person meant.

\begin{cousinweb}
\textbf{\gu{હરીફ}} (\emph{harīph}) is usually traced to Arabic \emph{ḥarīf}, \textquotedblleft{}a rival, a fellow player\textquotedblright{}, through Persian. The contest itself is \textbf{\gu{હરીફાઈ}}.
\end{cousinweb}

\begin{grammarlens}[title={What you've built}]
One more word for describing what you lived and felt.
\end{grammarlens}

\subsection*{Your turn}
\begin{itemize}
\raggedright
  \item \emph{Say it:} \emph{harīph}
  \item \emph{Say it:} \emph{harīph}, once more
  \item \emph{Say it:} \emph{āje pahelvānno harīph tenā kartā̃ moṭo che.} — Today the wrestler's rival is bigger than he is.
  \item \emph{Recall:} say the Gujarati for a wrestling ground, then the Gujarati for a wrestler, then say \emph{harīph} again
\end{itemize}

\begin{tcolorbox}[breakable,colback=teal!4,colframe=teal!35!black,title={Before you move on}]
What does \gu{હરીફ} mean? (\textquotedblleft{}A rival\textquotedblright{}.) Say it once more.
\end{tcolorbox}
\section[\texorpdfstring{\gu{પછાડવું} (pachāḍvũ)}{પછાડવું (pachāḍvũ)}]{\gu{પછાડવું} (pachāḍvũ) — to throw down}

\label{lesson:GU-C527-pachadvu}



\begin{quote}
Before the new one: say the Gujarati for a wrestler, then the Gujarati for a rival.
\end{quote}

\subsection*{You'll want to know: \gu{પછાડવું}}
\textbf{\gu{પછાડવું}} — \emph{pachāḍvũ} — \textquotedblleft{}to throw down\textquotedblright{}.

\begin{quote}
\gu{પહેલવાને} \gu{તેના} \gu{હરીફને} \gu{ધરતી} \gu{પર} \gu{પછાડ્યો}.
\end{quote}

\begin{quote}
\emph{pahelvāne tenā harīphne dhartī par pachāḍyo.}
\end{quote}

\begin{quote}
The wrestler threw his rival to the ground.
\end{quote}

Past, with \gu{તેણે} (he or she …): \gu{તેણે} \gu{હરીફને} \gu{પછાડ્યો}, \emph{teṇe harīphne pachāḍyo}.

\begin{cousinweb}
\textbf{\gu{પછાડવું}} (\emph{pachāḍvũ}) comes from \emph{*pracchaṭ}, a form CDIAL reconstructs, and CDIAL lists the Gujarati verb, \textquotedblleft{}to knock down\textquotedblright{} (CDIAL 8493).
\end{cousinweb}

\begin{grammarlens}[title={What you've built}]
One more word for describing what you lived and felt.
\end{grammarlens}

\subsection*{Your turn}
\begin{itemize}
\raggedright
  \item \emph{Say it:} \emph{pachāḍvũ}
  \item \emph{Say it:} \emph{pachāḍvũ}, once more
  \item \emph{Say it:} \emph{pahelvāne tenā harīphne dhartī par pachāḍyo.} — The wrestler threw his rival to the ground.
  \item \emph{Recall:} say the Gujarati for a wrestler, then the Gujarati for a rival, then say \emph{pachāḍvũ} again
\end{itemize}

\begin{tcolorbox}[breakable,colback=teal!4,colframe=teal!35!black,title={Before you move on}]
What does \gu{પછાડવું} mean? (\textquotedblleft{}To throw down\textquotedblright{}.) Say it once more.
\end{tcolorbox}
\section[\texorpdfstring{\gu{હંફાવવું} (hamphāvvũ)}{હંફાવવું (hamphāvvũ)}]{\gu{હંફાવવું} (hamphāvvũ) — to tire out (an opponent)}

\label{lesson:GU-C527-hamphavvu}



\begin{quote}
Before the new one: say the Gujarati for a rival, then the Gujarati for to throw down.
\end{quote}

\subsection*{You'll want to know: \gu{હંફાવવું}}
\textbf{\gu{હંફાવવું}} — \emph{hamphāvvũ} — \textquotedblleft{}to tire out (an opponent)\textquotedblright{}.

\begin{quote}
\gu{નાના} \gu{પહેલવાને} \gu{મોટા} \gu{હરીફને} \gu{ઘણો} \gu{હંફાવ્યો}.
\end{quote}

\begin{quote}
\emph{nānā pahelvāne moṭā harīphne ghaṇo hamphāvyo.}
\end{quote}

\begin{quote}
The small wrestler wore out his big rival.
\end{quote}

Past, with \gu{તેણે} (he or she …): \gu{તેણે} \gu{હરીફને} \gu{હંફાવ્યો}, \emph{teṇe harīphne hamphāvyo}.

\begin{cousinweb}
\textbf{\gu{હંફાવવું}} (\emph{hamphāvvũ}) is built from \textbf{\gu{હાંફવું}}, \textquotedblleft{}to pant\textquotedblright{}, with the causative \emph{-āv-}: to leave an opponent panting. CDIAL files the family under \emph{*hamph}, \textquotedblleft{}pant\textquotedblright{}, an imitative form (CDIAL 13973).
\end{cousinweb}

\begin{grammarlens}[title={What you've built}]
One more word for describing what you lived and felt.
\end{grammarlens}

\subsection*{Your turn}
\begin{itemize}
\raggedright
  \item \emph{Say it:} \emph{hamphāvvũ}
  \item \emph{Say it:} \emph{hamphāvvũ}, once more
  \item \emph{Say it:} \emph{nānā pahelvāne moṭā harīphne ghaṇo hamphāvyo.} — The small wrestler wore out his big rival.
  \item \emph{Recall:} say the Gujarati for a rival, then the Gujarati for to throw down, then say \emph{hamphāvvũ} again
\end{itemize}

\begin{tcolorbox}[breakable,colback=teal!4,colframe=teal!35!black,title={Before you move on}]
What does \gu{હંફાવવું} mean? (\textquotedblleft{}To tire out (an opponent)\textquotedblright{}.) Say it once more.
\end{tcolorbox}
